\documentclass[11pt]{article}

\usepackage[letterpaper,margin=0.7in]{geometry}
\usepackage{enumitem}
\usepackage{microtype}
\usepackage{titlesec}
\usepackage[hidelinks]{hyperref}

\pagestyle{empty}
\setlength{\parindent}{0pt}
\setlength{\parskip}{4pt}
\setlist[enumerate]{leftmargin=1.5em,itemsep=2pt,topsep=2pt}
\titleformat{\section}{\large\bfseries}{\thesection}{0.6em}{}
\titlespacing*{\section}{0pt}{8pt}{2pt}

\begin{document}

\begin{center}
  {\Large\bfseries APC 524 Project Proposal: hybopt}\\[2pt]
  {\large Embedding Trained Neural Networks in Mixed-Integer Optimization}\\[2pt]
  Brooke Soobrian \quad\textbullet\quad Harshit Verma \quad\textbullet\quad Zhe Li\\[2pt]
  Fall 2026
\end{center}

\section{Scope}

Artificial neural networks (ANNs) are widely used as surrogate models in engineering optimization.
Embedding a trained ANN in an optimization model, however, often makes the model computationally intractable because ANNs are nonconvex and complex.
Of the two most common ANN types, ReLU networks are typically formulated as mixed-integer linear programs (MILPs) through
a big-M reformulation, in which binary variables represent active and inactive neurons; however,
the big-M coefficients (bounds) can grow exponentially with network depth.
Tanh networks yield nonconvex nonlinear programs, whose tractability depends strongly on the reformulation of tanh.
The literature offers several solution methods for both cases, yet researchers have no common tool to implement and compare them.

We propose to build \textbf{hybopt}, a Python package based on Pyomo that embeds a trained ANN in an optimization model.
hybopt will support two use cases: (i) optimizing the ANN output directly, when the ANN is a surrogate for the entire problem,
and (ii) embedding the ANN as one constraint block inside a larger user-defined optimization model.
In both cases, the user chooses a solution method for each activation function.
We focus on dense feed-forward networks with bounded inputs, one hidden activation type, and a linear output layer;
network training and new solver development are out of scope.

\section{Content}

\begin{enumerate}
  \item \textbf{ReLU networks.} From Plate et al.~[1]: big-M with interval-arithmetic bounds, big-M with LP-based bound tightening (OBBT),
                                a-posteriori weight scaling, and scaling followed by bound tightening.
  \item \textbf{Exact tanh formulations.} From Schweidtmann and Mitsos~[2]: full-space formulations with alternative algebraic rewrites of tanh
                                          and a reduced-space formulation, solved to global optimality with SCIP.
  \item \textbf{Approximate tanh formulations} (our extension beyond [2]). A piecewise-linear (PWL) approximation of each tanh over its pre-activation bounds,
                                                with an SOS2 or binary encoding, reformulating the ANN as an MILP.
                                                A piecewise-quadratic (PWQ) approximation fits quadratic pieces separately on the convex and concave branches
                                                of tanh and reformulates the ANN as a mixed-integer quadratically constrained program (MIQCP).
                                                In both cases, the user controls the number of breakpoints or a target fit tolerance.
  \item \textbf{Benchmarks.} Trained ANN approximations of the Peaks, Ackley, and Himmelblau test functions, on which we compare the methods by solution time, model size, and optimality gap.
\end{enumerate}

\section{Implementation}

hybopt will build every formulation as a Pyomo model; Gurobi or HiGHS will solve the MILPs and SCIP will be used to solve the nonlinear models.
A common method interface will let users switch solution methods.
We will implement the ReLU methods first and then extend the interface to tanh.
The project is also an exercise in research-software engineering: we will use a GitHub repository with a pull-request and review workflow,
pytest suites that check each formulation against direct network evaluation and small problems with known solutions,
continuous integration through GitHub Actions, pre-commit quality checks, and documentation built with mkdocs alongside example notebooks.

\vspace{2pt}
{\small
\textbf{References}\\
{[1]} C.~Plate et al. An analysis of optimization problems involving ReLU neural networks. \emph{Optim.\ Eng.}, 2026. doi:\nolinkurl{10.1007/s11081-026-10075-8}\\
{[2]} A.~M. Schweidtmann and A.~Mitsos. Deterministic global optimization with artificial neural networks embedded. \emph{J.\ Optim.\ Theory Appl.}, 180:925--948, 2019. doi:\nolinkurl{10.1007/s10957-018-1396-0}
}

\end{document}
